\section{Выводы}
Созданы две динамические библиотеки для варианта 13 с единым контрактом вычисления производной и площади. Первая программа получает информацию о библиотеке при компоновке, вторая открывает модуль и ищет функции во время работы. Команда \texttt{0} переключает сразу обе реализации.

Подключение при компоновке упрощает клиентский код, но фиксирует зависимость от библиотеки до запуска. Явная загрузка даёт возможность выбирать реализацию в работающем приложении, однако переносит на него управление дескрипторами, адресами и ошибками загрузки. При обоих способах необходимо сохранять двоичную совместимость контракта.

Особое внимание уделено освобождению ресурсов: старый модуль закрывается только после получения новых адресов, а последующие вызовы используют новую библиотеку. Интеграционные проверки Windows подтвердили работу обеих программ, переключение и обработку ряда ошибок. Linux-ветка реализована, но для интегрированной версии ещё требует отдельной проверки. Производительность способов подключения не измерялась.
